% Chapter 1
\chapter{مقدمه}

\section{مقدمه}

در حال حاضر، استفاده از دوقلوهای مجازی در صنایع مختلف، از جمله تولید، بهداشت و درمان، حمل و نقل و حتی شهرسازی، در حال افزایش است. این تکنولوژی به سازمان‌ها این امکان را می‌دهد که فرآیندها و محصولات خود را به شکل مؤثرتری مدیریت کنند و تصمیمات هوشمندانه‌تری اتخاذ نمایند. به طور خاص، توانایی پیش‌بینی خرابی‌ها و مشکلات پیش‌رو از طریق تجزیه و تحلیل داده‌های جمع‌آوری شده، به بنگاه‌های اقتصادی کمک می‌کند تا به سمت کاهش هزینه‌های عملیاتی و افزایش بهره‌وری حرکت کنند.


تاریخچه استفاده از مدل‌های دیجیتال به سال‌ها قبل برمی‌گردد. مشهورترین مثال از این کاربردها، سیستم شبیه‌سازی که در طی ماموریت آپولوی 13 توسط ناسا ایجاد شد، است. در این ماموریت، زمانی که یکی از مخازن اکسیژن در فضاپیما منفجر شد، تیم کنترل پرواز در هیوستون با استفاده از یک مدل شبیه‌سازی شده از فضاپیما و اجزای آن، توانستند راه‌حلی برای عبور از این مشکل پیدا کنند. این واقعه نشان داد که چطور شبیه‌سازی‌ها می‌توانند به مهندسان کمک کنند تا از راه دور روندهای پیچیده را در محیطی امن آزمایش کنند و چالش‌های زندگی واقعی را شبیه‌سازی کنند. هرچند این مدل‌ها در آن زمان، به‌عنوان دوقلوهای مجازی شناخته نمی‌شدند، اما تأثیر آنها در درک عمیق‌تر از ارتباط میان فضاهای فیزیکی و مجازی غیرقابل انکار بود.
\cite{b3}

در نهایت، با پیشرفت مستمر فناوری‌های دیجیتال، عصر جدیدی از نوآوری و کارایی در انتظار صنایع متعدد است. دوقلوهای مجازی به عنوان یکی از مؤلفه‌های کلیدی این تحول، نقش مهمی در ارتباط و همگرایی میان دنیای فیزیکی و دیجیتال ایفا خواهند کرد. بدین ترتیب، می‌توان انتظار داشت که کاربردهای دوقلوهای مجازی در آینده نزدیک گسترش یابند و به بهبود مستمر کیفیت خدمات و محصولات منجر شوند.


\section{تعریف دوقلوی مجازی}


دوقلوهای مجازی 
\LTRfootnote{Digital Twins}
 به عنوان یک فناوری نوین در حوزه مهندسی و مدیریت سیستم‌ها، از جمله تحولات مؤثر در عصر فناوری اطلاعات محسوب می‌شود. این مفهوم بیشتر در زمینه مدیریت چرخه عمر محصولات
 \LTRfootnote{Product Lifecycle Management}
  مورد توجه قرار گرفت. 
  در واقع، دوقلوهای مجازی به معنای وجود یک مدل دیجیتال وابسته به یک شیء فیزیکی در دنیای واقعی هستند که از طریق تبادل داده‌ها، همگام‌سازی و تجزیه و تحلیل آنچه در دنیای واقعی اتفاق می‌افتد، امکان بهینه‌سازی عملیات و پیش‌بینی مشکلات را فراهم می‌سازند. یک دقلوی مجازی می‌تواند به روش‌های مختلفی با همتای حقیقی خود ارتباط برقرار کند و بر روی آن اثر گذاشته یا از آن اثر پذیرد.

دوقلوهای مجازی به عنوان یک پیوند مؤثر میان دنیای فیزیکی و دیجیتال شناخته می‌شوند که نه تنها نیاز به تبادل داده‌ها را تسهیل می‌کنند، بلکه بر پایه آن، امکان تعامل بلادرنگ و بهینه‌سازی فرآیندها را نیز ایجاد می‌کند.
\cite{b2}
 این مفهوم از طریق ایجاد یک اتصالات دائمی و یا حداقل دوره‌ای بین فضای واقعی و مجازی به وجود می‌آید که در آن، داده‌های دنیای واقعی به فضای مجازی منتقل می‌شوند و برعکس. 
از این رو، دوقلوهای مجازی می‌توانند حاوی اطلاعات دقیقی از وضعیت فعلی و پیش‌بینی‌های مربوط به شرایط آینده یک سیستم باشند.

می‌توان دوقلوی‌های مجازی را به سه بخش اصلی تقسیم کرد:
\cite{b2}
\begin{enumerate}
	\item 
	بخش اول، بخش مجازی که در یک فضای مجازی نیز قرار دارد.
	\item 
	بخش دوم، بخش واقعی و حقیقی که در یک فضای حقیقی و واقعی قرار دارد.
	\item 
	بخش سوم، جریان داده‌ها میان بخش واقعی و بخش حقیقی وجود دارد.
\end{enumerate}

در واقع دوقلوی مجازی به مجموعه‌ای از مدل‌‌ها و ماشین‌های فیزیکی یا مجازی گفته می‌شود که عملکرد و چرخه حیات یک موجودیت فیزیکی، را شبیه‌سازی\LTRfootnote{simulating}، تقلید\LTRfootnote{emulating} کرده یا آینه‌وار\LTRfootnote{mirroring} نمایش می‌دهند. البته این دوقلو تنها به موجودیت‌های فیزیکی خلاصه نمی‌شود و می‌توان در لایه‌های مختلفی دوقلو‌های مجازی را بررسی کرد:
\begin{itemize}
	\item 
	دوقلوی مجازی در سطح اشیاء که در آن سعی بر مجازی‌سازی و ایجاد دوقلوی مجازی از اشیاء و موجودیت‌های مختلف است.
	\item 
	دوقلوی مجازی در سطح فرایند که از مجموعه‌ای از کنش‌ و برهم‌ کنش موجودیت‌های مختلف ایجاد می‌شود.
	\item 
	دوقلوی مجازی سامانه که از مجموعه‌ای از فرایند‌ها ایجاد می‌شود و قابلیت تعریف و تفسیر سامانه‌ها را دارد.
\end{itemize}

علاوه‌بر انواع مختلفی از دوقلو‌های مجازی که به آن‌ها اشاره شد، می‌توان موارد دیگری را نیز در نظر گرفت که با ترکیب کردن آن‌ها، قدرت ایجاد دوقلو‌های مجازی بزرگ‌تر، جامع‌تر و کامل‌تری را به صورت قابل توجهی افزایش می‌دهد به طوری که بتوان دوقلو‌های مجازیی یک انسان یا ویژگی‌های آن‌را نیز ایجاد نمود.

دوقلوهای مجازی تنها یک مدل یا شبیه‌سازی ساده نیست، بلکه یک مدل زنده،‌هوشمند و در حال تکامل و توسعه است که به عنوان یک همتای مجازی برای موجودیت حقیقی ایجاد می‌شود.

\section{عامل}

به طور کلی برای 
\textit{عامل}
یک تعریف واحد و یکتا که در همه زمینه‌ها کاربرد داشته باشد، یافت نمی‌شود. این موضوع به این دلیل است که در زمینه‌های مختلف، می‌تواند ویژگی‌های مختلفی داشته باشد و همچنین اشکال فیزیکی و غیر فیزیکی متنوعی را می‌تواند شامل شود، از ربات‌ها گرفته تا شبکه‌های کامپیوتر و حتی عامل‌های نرم‌افزاری.
یکی از تعاریفی که در حوزه هوش‌مصنوعی ارائه شده به صورت زیر است:

\textit{
یک برنامه نرم افزاری است که می تواند با محیط خود تعامل داشته باشد، داده ها را جمع آوری کند و از داده ها برای انجام وظایف خود و تصمیم‌گیری در جهت رسیدن به اهداف از پیش تعیین شده استفاده کند.}\cite{b4}

در برخی‌ از منابع، ویژگی‌های زیر برای یک عامل در نظر گرفته شده است:\cite{b5}


\begin{description}
    \item[$\bullet$ قرارگیری\footnotemark:] \footnotetext{Situatedness}
 قرارگیری به معنای تعامل یک عامل با محیط از طریق حسگرها و عملگرها است. محیطی که یک عامل در آن حضور دارد، بخش جدایی‌ناپذیری از طراحی آن به شمار می‌رود و ورودی‌ها به‌طور مستقیم به‌عنوان نتیجه‌ی تعاملات عامل با محیط دریافت می‌شوند. این عامل‌ها به صورت مستقیم یا غیر مستقیم بر محیط تأثیر می‌گذارند و تاثیر می‌پذیرند.
    \item[$\bullet$ خودمختاری \footnotemark:] \footnotetext{Autonomy}
خودمختاری توانایی یک عامل را برای انتخاب عمل‌های خود به‌طور مستقل و بدون هیچ مداخله‌ای از سوی دیگر عامل‌ها یا انسان‌ها نشان می‌دهد. این ویژگی عامل را در برابر تأثیرات خارجی محافظت می‌کند و حالت‌های داخلی آن را از ناپایداری ناشی از اختلالات خارجی مجزا\LTRfootnote{Isolate} می‌سازد. بدین ترتیب، عامل قادر است به‌طور مؤثری در محیط خود عمل کند و تصمیم‌های بهتری اتخاذ نماید.
    \item[$\bullet$ توانایی استنتاج\footnotemark:] \footnotetext{Inferential Capability}
 توانایی استنتاج به کمترین توانایی عامل جهت نتیجه‌گیری و تصمیم‌گیری با استفاده از داده‌های موجود، گفته می‌شود. عامی می‌تواند با استنتاج مشاهده‌ها از طریق تعمیم اطلاعات، به میزانی از درک وضعیت‌های مختلف برسد و بر اساس آن بهترین تصمیم را بگیرد.
    \item[$\bullet$ واکنش‌دهی\footnotemark:] \footnotetext{Responsiveness}
  واکنش‌دهی به معنای توانایی درک شرایط محیط و پاسخ به آن به‌طور به‌موقع است. این ویژگی برای بسیاری از برنامه‌های بی‌درنگ\LTRfootnote{Realtime} اهمیت بالایی دارد، زیرا عامل باید قادر باشد بلافاصله به تغییرات محیطی واکنش نشان دهد و ضمن پایش مداوم شرایط، تصمیمات مناسبی را اتخاذ کند.
    \item[$\bullet$ پیش‌فعالی\footnotemark:] \footnotetext{Pro-activeness}
    
 پیش‌فعالی به معنای توانایی عامل در پاسخ به رفتارهای فرصت‌طلبانه و گرفتن تصمیمات بهینه و هدفمند است. عامل باید بتواند به صورت فعال به فرصت‌ها و چالش‌های جدید پاسخ دهد و در محیط پویا و متغیر به‌صورت موثرتر عمل نماید.
    \item[$\bullet$ رفتار اجتماعی \footnotemark:] \footnotetext{Social Behaviour}
رفتار اجتماعی به قابلیت تعامل عامل با منابع خارجی در مواقع ضروری اشاره دارد تا به اهداف خاص دست یابد. علی‌رغم اینکه تصمیم‌گیری عامل باید آزادانه و بدون مداخله خارجی باشد، باید بتواند به‌طور مؤثر با دیگر عامل‌ها و نهادهای انسانی ارتباط برقرار کند و دانش خود را به اشتراک بگذارد. این امر به عامل کمک می‌کند تا از تجربیات دیگران بیاموزد و به حل مسائل مشترک کمک کند، که خود به بهبود همکاری و کارایی کلی سیستم‌های چند عاملی\LTRfootnote{MAS} منجر می‌شود.
\end{description}
علاوه‌بر ویژگی‌هایی که بیان شد، عامل‌ها می‌توانند بسته به کارکرد خود، ویژگی‌های دیگری نیز داشته باشند. همچنین بنابر عملکری که هر عامل باید داشته باشد، ممکن است برخی از ویزگی‌های بیان شده را دارا نباشد. پس هر عامل را می‌توان بر اساس این ویژگی‌ها سنجید و بررسی کرد.





\section{تعریف سیستم‌های چند ‌عاملی}
یک سیستم چندعاملی\LTRfootnote{Multi Agent System}
از تعدادی عامل ایجاد می‌شود که به صورت مستقل یا گروهی، با همکاری یا به صورت رقابتی، در یک محیط برای رسیدن به یک هدف یا اهداف مشخص، تشکیل شده است. این عامل‌ها می‌توانند دانش خود را با هم به اشتراک بگذارند و به کمک یکدیگر نتیجه و تصمیم بگیرند و یا هرکدام به صورت مستقل عمل کرده و بر روی محیط تاثیر گذارد.

سیستم‌های چند عاملی به دلیل قابلیت‌های منحصربه‌فردی که ارائه می‌دهد، قابل استفاده در زمینه‌ها و کاربرد‌های مختلف است. چند مورد از قابلیت‌های این سیستم را در اینجا بررسی می‌کنیم:\cite{b5}

\begin{itemize}
	\item
	انجام محاسبات و عمل‌ها به صورت موازی و غیر همزمان باعث افزایش سرعت و کارایی می‌شود سیستم می‌شود.
	\item 
	هر زمان که یک یا چند عامل از سامانه خارج شود و دچار مشکل شود، کار کل سامانه مختل نمی‌شود و به کار خود ادامه می‌دهد. به این صورت میزان اطمینان‌پذیری و پایداری سامانه را افزایش می‌یابد.
	
	\item
	در صورتی که نیاز باشد می‌توان عامل‌های تازه به سامانه اضافه کرد که این مورد انعطاف‌پذیر و مقیاس‌پذیر سامانه را فراهم می‌کند.
	
	\item 
	هزینه مربوط به هر عامل می‌تواند کمتر از سامانه‌های مرکزی و غیر توزیع‌شده باشد.
	\item 
	توسعه عامل‌ها به دلیل کوچک‌تر بودن و مشخص بودن کار آن‌ها، راحت‌تر و ساده‌تر و با هزینه کم‌تر انجام‌پذیر است.
	
\end{itemize}

در سیستم‌‌های چند عاملی می‌توان مسئله‌ها و وظیفه اصلی و بزرگ را به مسئله‌ها و وظیفه‌های کوچک‌تر تقسیم کرد و انجام هرکدام از آن‌ها را به عامل‌های جداگانه سپرد. با این کار می‌توان کار‌های بزرگی که انجام آن‌ها توسط یک عامل بسیار سخت و پیچیده  و هزینه‌بر است را به راحتی انجام داد.\cite{b6}

\section{هوش مصنوعی در دوقلو‌ی دیجیتال}

امروزه کمتر حوزه‌ای می‌توان یافت که در آن هوش مصنوعی یافت نشود. امروزه در زمینه‌های مختلف استفاده از روش‌ها و الگوریتم‌های هوش مصنوعی کاربرد فراوانی یاف	ته است. تعاریف مختلف و متنوع‌ای از هوش مصنوعی وجود دارد که یکی از آن‌ تعاریف به شرح زیر است:

\textit{
	هوش مصنوعی به علم و مهندسی ساخت ماشین های هوشمند، به ویژه برنامه های کامپیوتری هوشمند گویند به طوری که توانایی کسب، فهمیدن و بکارگیری دانش را داشته‌باشد. همجنین به استفاده از رایانه برای درک هوش انسانی مربوط می شود.
	}\cite{b13}

هوش مصنوعی در اوایل ظهور خود، با اینترنت اشیاء\LTRfootnote{Internet of things (IoT)}
پیوند خورد و محصولات تازه‌ای وارد بازار کرد. تعریف هوش مصنوعی اشیاء\LTRfootnote{Artificial Intelligence of things (AIoT)}
از ترکیب فناوری هوش مصنوعی و اینترنت اشیاء ایجاد شده‌است. در گذشته دوقلوی دیجیتال و هوش مصنوعی دو مفهوم کاملا مجزایی دانسته می‌شدند اما امروزه با گسترش و تنوع چشم‌گیر حسگر‌ها و همچنین پررنگ شدن اهمیت تحلیل داده‌ها و پاسخگویی و اقدام هوشمندانه و سریع بر اساس داده، قابل جدا سازی نیستند. هوش مصنوعی به عنوان بخشی از لایه‌های دوقلوی دیجیتال کاملا پذیرفته شده و در حال استفاده است. این موضوع را می‌توان در مثال‌ها و نمونه‌هایی که از دوقلو‌ی دیجیتال آورده‌ خواهد شد به وضوح، مشاهده کرد.

دوقلوی‌ دیجیتال با هوش مصنوعی از داده‌های فرآیند‌ها برای پیش‌بینی اینکه فرآیند در آینده چگونه خواهد بود، استفاده می‌کند. این عملکرد نسبت به گذشته دوقلوی دیجیتال که تنها یک شبیه‌سازی و نمایش سه‌بعدی از فرایند بود، بسیار کاربردی‌تر و مهم‌تر است.


\section{خلاصه‌ای از کار‌ها و نتایج پیشین}


بر اساس مقاله 
\cite{b7}
دوقلو‌های مجازی معمولا از سه جزو ایجاد می‌شود. فضای واقعی دارای شئ واقعی، فضای مجازی دارای شئ مجازی و ارتباط میان شئ مجازی و واقعی و انتقال داده ‌میان آن‌ها.

طبق ادعایی که در مقاله‌های 
\cite{b8}
و
\cite{b9}
صورت گرفته‌است، دوقلو‌های مجازی تنها یک شبیه‌سازی ساده نیست. بلکه یک مدل زنده، هوشمند و در حال تکامل است. دوقلوی مجازی، چرخه عمر همتای حقیقی خود را دنبال می‌کند تا بتواند آن را نظارت\LTRfootnote{Monitor}
و کنترل\LTRfootnote{Control}
کرده و همچنین کارکرد آن را بهبود بخشد. دوقلوی مجازی به صورت مستمر با فضا و محیط واقعی در ارتباط بوده و به صورت بی‌درنگ با آن، داده‌ ردوبدل می‌کند.


مقاله 
\cite{b10}
به بررسی دوقلو‌های مجازی در اینترنت وسایل‌ نقلیه\LTRfootnote{Internet of Vehicles (IoF)}
پرداخته است. در واقع این مقاله، یک سکو\LTRfootnote{Platform}
جهت به اشتراک گذاری داده، بین وسایل نقلیه موجود در سطح شهر را معرفی می‌کند. در این مقاله، یک دوقلوی دیجیتال، به صورت یک بازتاب از وسیله نقلیه در فضای ابری تعریف می‌شود که داده‌ها را با وسیله نقلیه در هر لحظه همگام می‌کند. داشتن این دوقلوی مجازی در بستر ابری، باعث می‌شود که بتوان با به اشتراک‌گذاری داده‌ها بدون محدودیت فاصله و پهنای باند، کارایی رانندگی را افزایش داده و تصمیمات دقیق‌تری بر اساس داده‌ها گرفته‌ شود.\\
در نهایت مقاله 
\cite{b11}
یک چارچوب\LTRfootnote{Framework}
به نام 
\lr{MADTwin}\LTRfootnote{Multi Agent Digital Twin}
ارائه می‌دهد که برای مدل‌سازی دوقلوی دیجیتال برپایه عامل‌\LTRfootnote{Agent}
مناسب است. از این چارچوب برای هوشمند‌سازی یک انبار‌ توسط ربات‌ها استفاده کرده تا بتوان به طور بهینه هر ربات را به سوی اهداف محلی خود و اهداف کل سیستم به صورت مشارکتی یا مستقل هدایت کند.





% \begin{figure}[H]
%   \begin{subfigure}{0.5\textwidth}
%     \centering
%     \includegraphics[width=\linewidth]{low00081.png}
%     \caption{تصویر کم نور}
%     \label{fig:subfig1}
%   \end{subfigure}%
%   \begin{subfigure}{0.5\textwidth}
%     \centering
%     \includegraphics[width=\linewidth]{normal00081.png}
%     \caption{تصویر بهبودیافته}
%     \label{fig:subfig2}
%   \end{subfigure}
%   \caption{نمایش تصویر کم نور و نمونه بهبودیافته آن\cite{b27}}
%   \label{fig:whole}
% \end{figure}

% \section{پیامدهای بصری تصویربرداری در نور کم}
 
% تصویربرداری در نور کم پیامدهای بصری خاصی را به همراه دارد که می‌تواند بر کیفیت تصاویر تأثیر گذار باشد. در ادامه به برخی از این پیامدها اشاره می‌شود که نیاز به راهکارهای مناسب برای بهبود تصاویر در این شرایط را تاکید می‌کند.

% \subsection{افزایش نویز}

% یکی از بارزترین پیامدهای بصری تصویربرداری در نور کم، افزایش قابل توجه نویز \LTRfootnote{Noise} تصویر است. نویز به صورت تغییرات ناخواسته در مقادیر پیکسل\LTRfootnote{Pixel}‌ها ظاهر می‌شود. نویز به ویژه در مناطقی از تصویر که در حالت ایده‌آل باید یکنواخت باشد مشخص است. کمبود فوتون\LTRfootnote{Photon}‌های ورودی در طول دریافت نور کم باعث کاهش نسبت سیگنال به نویز\LTRfootnote{Signal-to-noise Ratio} می‌شود. در نتیجه تصاویر سطوح بالائی از نویز را نشان می‌دهند که نه تنها کیفیت زیبائی‌شناختی تصویر را به خطر می‌اندازد، بلکه مانع تحلیل و تفسیر دقیق صحنه می‌شود. این شکل از نویز اغلب به‌عنوان نویز روشنائی\LTRfootnote{Luminance Noise} شناخته می‌شود.

% \subsection{کاهش تضاد درخشندگی}

% شرایط کم نور اغلب منجر به کاهش تضاد درخشندگی \LTRfootnote{Contrast} بین عناصر مختلف در تصویر می‌شود. تضاد درخشندگی یک جنبه اساسی از کیفیت تصویر است که نشان‌دهنده تفاوت در سطح روشنائی روشن‌ترین و تاریک‌ترین مناطق تصویر است. در شرایط کم نور این تفاوت کاهش می‌یابد و باعث می‌شود اشیا، لبه‌ها و جزئیات کمتر متمایز به نظر برسند. کاهش تضاد درخشندگی توانائی تشخیص و شناسائی اشیاء، به‌ویژه آن‌هایی که در مناطق سایه‌دار یا کم‌نور قرار دارند را مختل کند.

% \subsection{از دست دادن جزئیات تصویر}

% روشنائی ناکافی در شرایط نور کم می‌تواند منجر به از دست دادن جزئیات مهم تصویر شود. بافت‌های ظریف، ویژگی‌های دقیق و ساختارهای پیچیده ممکن است نامشخص یا کاملا مبهم شوند. این از دست دادن جزئیات پیامدهایی برای کاربردهای مختلف از جمله تشخیص اشیاء، تصویربرداری پزشکی و نظارت دارد.

% \section{پیامدهای مفهومی تصویربرداری در نور کم}

% تصاویری که در شرایط کم نور ثبت می‌شوند معمولا از کیفیت زیبائی شناختی و انتقال نامطلوب اطلاعات رنج می‌برند. تصویربرداری در نور کم بر ویژگی‌های بصری و کاربردهای عملی آن‌ها تاثیر می‌گذارد، در قسمت قبلی در مورد پیامدهای بصری بحث شد. پیامدهای مفهومی منجر به ارسال پیام اشتباه از طریق تصاویر می‌شود. به‌عنوان مثال در کاربردهایی مانند تشخیص اشیا یا چهره\LTRfootnote{Object or Face Detection}، اطلاعات نادرست یا ناقص در تصاویر کم نور می‌تواند مانع از شناسائی دقیق اشیاء یا افراد شود.

% \begin{figure}
%   \begin{subfigure}{0.5\textwidth}
%     \centering
%     \includegraphics[width=\linewidth]{r0ea0825ft.png}
%     \caption{تصویر کم نور}
%     \label{fig:subfig1}
%   \end{subfigure}%
%   \begin{subfigure}{0.5\textwidth}
%     \centering
%     \includegraphics[width=\linewidth]{r0ea0825ft_.png}
%     \caption{تصویر بهبودیافته}
%     \label{fig:subfig2}
%   \end{subfigure}
%   \caption{نمایش پیامدهای بصری و مفهومی در تصویر کم نور\cite{b27}}
%   \label{fig:whole}
% \end{figure}


% \section{پیشرفت‌های تکنولوژی و چالش‌های پایدار}

% تا به امروز فناوری تصویربرداری، حساسیت حسگرها و الگوریتم‌های کاهش نویز با سرعت بی‌سابقه‌ای تکامل یافته است. این پیشرفت‌ها امکانات تصویربرداری در نور کم را گسترش داده و افق‌های جدیدی را برای کاربردهای گوناگون در حوزه‌های مختلف باز کرده است. با این حال لازم است اذعان شود که تصویربرداری در نور کم کاری پیچیده است. شرایط محیطی نامطلوب مانند تاریکی شدید، مه یا آب و هوای نامساعد می‌تواند پیشرفته‌ترین تجهیزات را تحت فشار قرار دهند. محدودیت‌های فنی نظیر محدودیت‌های فناوری حسگر یا قابلیت‌های لنز می‌توانند چالش‌های پیش رو در شرایط کم نور را تشدید کنند. لازم به ذکر است که هزینه بالای دوربین‌های تخصصی، لنزها و لوازم جانبی که قادر به ارائه عملکرد استثنایی در نور کم هستند، می‌تواند دسترسی به قابلیت‌های تصویربرداری در شرایط کم نور را محدود می‌کند. این امر اهمیت ارتقاء کیفیت این‌گونه تصاویر از طریق روش‌های بهبود کیفیت تصاویر کم را آشکار می‌کند.

% \section{اهمیت بهبود کیفیت تصاویر کم نور}

% تلاش برای بهبود کیفیت تصاویر کم نور فراتر از ملاحظات زیبائی‌شناختی است چرا که تصاویر به‌عنوان یک عنصر اساسی برای اجرای موفقیت آمیز وظایف گوناگون بینائی ماشین عمل می‌کند. در قلمرو بینائی ماشین کیفیت داده‌های ورودی بسیار مهم است. تصاویر با نور کم که با نویز، کاهش تضاد درخشندگی و از دست دادن احتمالی جزئیات تغییر یافته اند، چالش‌های ذاتی برای تفسیر دقیق اطلاعات بصری ایجاد می‌کنند.

% این امر در کاربردهائی مانند تشخیص اشیاء، جائی که شناسائی و مکان‌یابی اشیاء در یک تصویر بسیار مهم است اهمیت می‌یابد و می‌تواند منجر به مکان‌یابی نادرست اشیاء شود. تقسیم‌بندی اشیاء\LTRfootnote{Object Segmentation} یکی دیگر از وظایف بینائی ماشین، به شدت بر توانائی تمایز دقیق اشیاء از پس‌زمینه آن‌ها متکی است. تصاویر با نور کم ممکن است مرزهای اشیاء را مبهم کنند و بخش‌بندی را مستعد خطا کنند. به‌طور مشابه وظایف تشخیص تصویر مانند شناسائی الگوها یا اشیاء درون تصاویر، به وجود ویژگی‌های واضح و قابل تشخیص بستگی دارد. شرایط کم نور می‌تواند قابلیت اطمینان الگوریتم‌های تشخیص را به خطر بیاندازد و منجر به تقسیم‌بندی‌های نادرست یا تلاش‌های ناموفق برای تشخیص شود. علاوه بر این اهمیت بهبود کیفیت تصاویر کم نور به کاربردهای حیاتی پزشکی گشترش می‌یابد. شرایط کم نور در زمینه‌هایی مانند رادیولوژی\LTRfootnote{Radiology}، آندوسکوپی و تشخیص پزشکی، جائی که تصویربرداری دقیق ضروری است، می‌تواند مانع از تشخیص ناهنجاری‌های ظریف یا آسیب‌شناسی شود.

% به صورت کلی با توجه به نقش محوری کیفیت تصویر در وظایف بینائی ماشین، پیش پردازش تصاویر در شرایط کم نور اهمیت حیاتی دارد. روش‌های پیش‌پردازش از جمله کاهش نویز، افزایش تضاد درخشندگی و بازیابی جزئیات در آماده‌سازی تصاویر در نور کم برای تجزیه و تحلیل بسیار مفید هستند. هدف این روش‌ها اصلاح پیامدهای بصری شرایط کم نور است و تضمین می‌کند که فرآیندهای بینائی ماشین بعدی بر اساس داده های تصویری با کیفیت بالا و با جزئیات است.

% % \begin{figure}[t]
% % \centering
% %\includegraphics[height=8cm,width=10cm]{wt.png}
% %\caption{ توربین بادی 12 کیلو‌وات (ساخته‌شده به وسیله‌ی چارلز فرانسیس براش) \cite{b1}.}
% %\label{wt}
% %\centering
% %\end{figure}



% %\section{اهداف و دستاوردهای تحقیق}


% %\section{ساختار گزارش}

% \section{نتیجه‌گیری}

% تصویربرداری در نور کم چالش‌های منحصر به فردی را ایجاد می‌کند که منجر به کاهش کیفیت و تفسیرپذیری تصاویر می‌شود. این چالش‌ها شامل افزایش نویز، کاهش تضاد درخشندگی و از دست دادن جزئیات می‌شود. این پیامدهای بصری می‌تواند برای کاربردهای مختلف، از جمله تشخیص اشیاء، تصویربرداری پزشکی و نظارت، مشکل‌ساز باشد.
% پیشرفت‌های تکنولوژیکی در زمینه تصویربرداری، حساسیت حسگرها و الگوریتم‌های کاهش نویز را بهبود بخشیده است. با این حال، تصویربرداری در نور کم همچنان یک کار چالش برانگیز است و نیاز به توسعه روش‌های جدید برای بهبود کیفیت تصاویر کم نور دارد.
% روش‌های پیش‌پردازش تصاویر، مانند کاهش نویز، افزایش تضاد درخشندگی و بازیابی جزئیات، می‌توانند برای اصلاح پیامدهای بصری شرایط کم نور استفاده شوند. این روش‌ها می‌توانند به بهبود کیفیت تصاویر کم نور و افزایش قابلیت اطمینان کاربردهای بینایی ماشین کمک کنند.
